\documentclass[12pt]{article}
\usepackage[utf8]{inputenc}
\usepackage[T1]{fontenc}
\usepackage{hyperref}
\title{Krevinian Language Resources for K-12 Curriculum Development}
\author{Compiled by Claude Sonnet 5.5 (claude-sonnet-5-5) for Andr\'as Kornai}
\date{}
\begin{document}
\maketitle

\section{Introduction}
Krevinian (also Kreewinian, Krewinisch; ISO 639-3 zkv) is an extinct Finnic language of Latvia. According to the sources found, it was spoken around Bauska in Zemgale (Semigallia) and became extinct by the middle of the nineteenth century. Its speakers were descended from Votes taken from Ingria in the 1440s. Krevinian is often treated as a variety of Votic. Very little material survives. The only substantial early documentation found in the searches is that of Johan Andreas Sj\"ogren in the mid-nineteenth century. No school materials, textbooks or revitalization programmes were found. A curriculum can therefore be built only on the surviving word lists and texts, and on the related Votic and Finnic languages.

\section{Primary Educational Resources}
No primers, textbooks or school programmes for Krevinian were found in the search results.

\section{Linguistic and Reference Materials}
\begin{itemize}
\item Wikipedia, ``Krevinian dialect'': overview of the variety and its documentation. \url{https://en.wikipedia.org/wiki/Krevinian_dialect}
\item Wikipedia, ``Kreevins'': history of the Krevin people of Zemgale. \url{https://en.wikipedia.org/wiki/Kreevins}
\item Ugri.net, ``Other Finnic languages'': overview of minor Finnic languages, which may include Krevinian. \url{https://www.ugri.net/in-english/uralic-languages/finnic-languages/other-finnic-languages}
\item LMU Munich e-learning, Finnic languages, part 6 (PDF): university teaching material on Finnic. \url{https://www.infuse.finnougristik.uni-muenchen.de/e-learning/finnic/o6_finnic_part6.pdf}
\item Linguistica Uralica (2025), a study of Krevinian place-name etymologies, which concludes that most proposals are doubtful (PDF). \url{https://kirj.ee/wp-content/plugins/kirj/pub/ling-2025-1-1-29_20250310122410.pdf}
\end{itemize}

\section{Religious and Community Resources}
No religious texts or community organizations specific to Krevinian were found in the searches.

\section{Digital and Audio Resources}
No audio, video or digital corpora were found. The extinct language has no recordings.

\section{Curriculum Development Context}
The language has no native speakers, so any programme would be a heritage-awareness and reconstruction project. Possible components are local history of Bauska and Zemgale, the story of the Votic captives, and the nineteenth-century word lists. Related living Finnic languages could supply comparative material. Little exists for the language itself, and the curriculum should say so honestly.
\end{document}